% ======================================================================
% Shared preamble -- Designing Agentic AI Solutions
% Loaded by main.tex and by each standalone chapter build.
% ======================================================================

% ---------- where to look for \input files ---------------------------
% Works both from the book folder (main.tex) and from a chapter folder
% (main-ch01.tex etc.), so chapters can say \input{ch01a} and
% \input{fig-integrated} without folder names.
\makeatletter
\@ifundefined{input@path}{\def\input@path{}}{}
\g@addto@macro\input@path{{./}{common/}{../common/}%
  {ch01/}{ch02/}{ch03/}{ch04/}{ch05/}{ch06/}{ch07/}%
  {ch08/}{ch09/}{ch10/}{ch11/}{ch12/}{ch13/}{appendix/}}
\makeatother

% ---------- fonts and page ------------------------------------------
\usepackage[T1]{fontenc}
\usepackage{lmodern}
\usepackage{microtype}
\usepackage[margin=1in]{geometry}

% ---------- maths, tables, lists, graphics ---------------------------
\usepackage{amsmath}
\usepackage{amssymb}% \checkmark
\usepackage{textcomp}
\usepackage{graphicx}
\graphicspath{{./}{../}{ch01/}{ch02/}{ch03/}{ch04/}{ch05/}{ch06/}{ch07/}%
  {ch08/}{ch09/}{ch10/}{ch11/}{ch12/}{ch13/}}
\usepackage[dvipsnames,table]{xcolor}
\usepackage{array}
\usepackage{booktabs}
\usepackage{tabularx}
\usepackage{enumitem}
\usepackage{multicol}
\usepackage{pdfpages}
% ---------- TikZ and the solution-graph notation ---------------------
\usepackage{tikz}
\usetikzlibrary{positioning,shapes.geometric,shapes.multipart,shapes.misc,
  arrows.meta,fit,backgrounds,calc}

% Building-block colours (Figure 1.1 and the integrated figure).
% Same palette as the solution-graph notation in sg-styles.tex.
\definecolor{agentfill}{HTML}{D6E4F7}
\definecolor{codefill}{HTML}{ECECEC}
\definecolor{mcpfill}{HTML}{D9F0E6}
\definecolor{ragfill}{HTML}{FBE8C4}
\definecolor{humanfill}{HTML}{F9DCE6}

% Solution-graph styles (sg agent, sg mcp, ...) and \sgtag
% ----------------------------------------------------------------------
% Simplified solution-graph styles (replace the old \tikzset{sg ...} block
% in common/preamble.tex with this file, or \input it after the old block).
% Three shapes only: box, cylinder, double-bordered box. Colour gives the type.
% Style names are unchanged, so every figure in the book picks them up.
% Needs TikZ libraries: shapes.geometric, shapes.multipart, arrows.meta, fit.
% ----------------------------------------------------------------------
\definecolor{sgagentc}{HTML}{D6E4F7}
\definecolor{sgsubc}{HTML}{EBF2FB}
\definecolor{sgcodec}{HTML}{ECECEC}
\definecolor{sgmcpc}{HTML}{D9F0E6}
\definecolor{sgstorec}{HTML}{FBE8C4}
\definecolor{sghumanc}{HTML}{F9DCE6}
\definecolor{sgguardc}{HTML}{FAD9C2}
\definecolor{sgcallc}{HTML}{1F4E8C}
\definecolor{sgtrustc}{HTML}{A33A3A}
\tikzset{
  sg base/.style={draw, thin, rounded corners=2pt, align=center, font=\footnotesize,
                  text width=2.1cm, minimum height=9mm, inner sep=2pt},
  % nodes
  sg trigger/.style={sg base, fill=white},
  sg agent/.style={sg base, fill=sgagentc},
  sg subagent/.style={sg base, fill=sgsubc},
  sg code/.style={sg base, fill=sgcodec},
  sg guard/.style={sg base, fill=sgguardc},
  sg mcp/.style={sg base, fill=sgmcpc, rectangle split, rectangle split parts=2,
                 rectangle split draw splits=false},
  sg human/.style={sg base, fill=sghumanc, double, double distance=1.2pt},
  sg store/.style={sg base, fill=sgstorec, cylinder, shape border rotate=90, aspect=0.18,
                   rounded corners=0pt},
  sg sor/.style={sg store, fill=white, dashed},
  % edges
  sg flow/.style={-{Stealth[length=2mm]}, thin},
  sg call/.style={-{Stealth[length=2mm]}, thin, sgcallc},
  sg data/.style={-{Stealth[length=2mm]}, thin, dashed},
  sg event/.style={-{Stealth[length=2mm]}, thick, dotted},
  sg handoff/.style={-{Implies}, line width=0.5pt, double, double distance=1.5pt},
  sg lab/.style={font=\scriptsize, fill=white, inner sep=1pt},
  % boundaries and tags
  sg boundary/.style={draw, gray, dashed, rounded corners=3pt},
  sg trust/.style={draw, sgtrustc, dash dot, rounded corners=3pt},
  sg btitle/.style={font=\scriptsize, anchor=south west, inner sep=1pt},
  sg tag/.style={fill=black, text=white, font=\tiny\sffamily\bfseries, inner sep=1.2pt, rounded corners=1pt},
}
\providecommand{\sgtag}{}
\renewcommand{\sgtag}[2]{\node[sg tag, anchor=south east] at (#1.north east) {#2};}


% ---------- code listings --------------------------------------------
\usepackage{listings}
\lstset{
  basicstyle=\ttfamily\small,
  keywordstyle=\color{RoyalBlue}\bfseries,
  commentstyle=\color{Gray}\itshape,
  stringstyle=\color{BrickRed},
  numbers=left, numberstyle=\tiny\color{Gray}, numbersep=6pt,
  frame=lines, framesep=4pt,
  breaklines=true, breakatwhitespace=true,
  showstringspaces=false, columns=fullflexible, keepspaces=true,
  captionpos=t, xleftmargin=1.5em, tabsize=4,
  upquote=true,
  % characters that listings cannot read directly from UTF-8 files
  extendedchars=true,
  literate={§}{{\S}}1 {é}{{\'e}}1 {è}{{\`e}}1 {ü}{{\"u}}1 {ö}{{\"o}}1
           {ä}{{\"a}}1 {€}{{\texteuro}}1 {£}{{\pounds}}1 {–}{{--}}1 {—}{{---}}1
           {’}{{'}}1 {‘}{{`}}1 {“}{{``}}1 {”}{{''}}1 {→}{{$\rightarrow$}}1
           {≥}{{$\geq$}}1 {≤}{{$\leq$}}1 {×}{{$\times$}}1 {±}{{$\pm$}}1,
}
\renewcommand{\lstlistingname}{Listing}

% ---------- boxes ----------------------------------------------------
\usepackage[most]{tcolorbox}
\definecolor{boxframe}{HTML}{3C3489}
\definecolor{boxfill}{HTML}{F3F2FB}
\definecolor{notefill}{HTML}{EEF6F2}
\definecolor{noteframe}{HTML}{0F6E56}
\definecolor{warnfill}{HTML}{FBEFEA}
\definecolor{warnframe}{HTML}{993C1D}

\tcbset{bookbox/.style={enhanced, breakable, boxrule=0.5pt, arc=2pt,
  left=6pt, right=6pt, top=4pt, bottom=4pt, fonttitle=\bfseries\sffamily}}

% Learning objectives at the start of each chapter
\newtcolorbox{objectives}{bookbox, colback=boxfill, colframe=boxframe,
  title={Learning objectives}}

% Definition of a key term
\newtcolorbox{definition}{bookbox, colback=white, colframe=boxframe,
  before upper={\textbf{\sffamily Definition.}\ }}

% A design principle, with an optional name: \begin{principle}[Name]
\newtcolorbox{principle}[1][]{bookbox, colback=boxfill, colframe=boxframe,
  before upper={\textbf{\sffamily Principle\ifx\relax#1\relax\else\ (#1)\fi.}\ }}

% Design note
\newtcolorbox{designnote}{bookbox, colback=notefill, colframe=noteframe,
  before upper={\textbf{\sffamily Design note.}\ }}

% Version note (framework and service versions)
\newtcolorbox{versionnote}{bookbox, colback=white, colframe=Gray, fontupper=\small,
  before upper={\textbf{\sffamily Version note.}\ }}

% Pattern card: \begin{patterncard}{Name}
\newtcolorbox{patterncard}[1]{bookbox, colback=notefill, colframe=noteframe,
  title={Pattern: #1}}

% Anti-pattern: \begin{antipattern}{Name}
\newtcolorbox{antipattern}[1]{bookbox, colback=warnfill, colframe=warnframe,
  title={Anti-pattern: #1}}
% Case at a glance (opening table of each case chapter)
\newtcolorbox{caseglance}{bookbox, colback=boxfill, colframe=boxframe,
  title={Case at a glance}}
% (The template environment is defined in appendix/appC.tex.)

% ---------- small commands -------------------------------------------
\newcommand{\keyterm}[1]{\textbf{#1}}

% Allow a line break after an underscore (long identifiers in \texttt)
\makeatletter
\DeclareRobustCommand{\_}{\ifmmode\nfss@text{\textunderscore}\else\textunderscore\allowbreak\fi}
\makeatother

% Inline code that may contain underscores: \code{get_order_status(order_number)}
% Long names may break after _ , . or ( so they fit narrow table columns.
\ExplSyntaxOn
\NewDocumentCommand{\code}{m}
  {
    \tl_set:Nx \l_tmpa_tl { \detokenize{#1} }
    \regex_replace_all:nnN { ([\_\,\.\(]) } { \1 \c{allowbreak} } \l_tmpa_tl
    \texttt{ \l_tmpa_tl }
  }
\ExplSyntaxOff

% Give a label to a chapter that is not written yet, so references print:
% \fakelabel{ch:business}{3}
\makeatletter
\newcommand{\fakelabel}[2]{\expandafter\gdef\csname r@#1\endcsname{{#2}{}{}{}{}}}
\makeatother

% ---------- citations and links (hyperref last) ----------------------
\usepackage[round]{natbib}
\usepackage[colorlinks=true, linkcolor=boxframe, citecolor=noteframe,
  urlcolor=boxframe]{hyperref}
